\section{命题分析}
\subsection{命题题目}
\ChallengeTitle

\par\noindent 命题企业：\ChallengeCompany{}。\cite{challenge}

\subsection{命题内容及答题要求}
\textbf{命题内容}\par
\ChallengeOriginalText

\par\textbf{答题要求}\par
\RequirementsOriginalText

\subsection{命题分析}
题面要求覆盖计算硬件、系统软件、智能任务和外部连接。结合现有资源，本方案将基础推理作为首期验证目标：先证明模型能够经过检查与转换，在板端获得可核对的计算结果，再扩展输入设备和外部服务。各项要求与设计工作的对应关系如下。

\begin{table}[H]
\centering
\caption{命题要求与当前方案的对应关系}
\small
\begin{tabularx}{\linewidth}{@{}P{.19\linewidth}YY@{}}
\toprule
要求类别 & 当前设计响应 & 验证或确认依据 \\
\midrule
处理器与硬件集成 &
以\ProcessorCore{}、\VectorUnit{}和\AcceleratorName{}构建 FPGA 计算原型，连接存储与外设 &
统一硬件配置、功能仿真、资源时序报告和板级运行记录 \\
开源操作系统 &
分别适配\RealTimeSystem{}与\LinuxDistribution{}，使用独立镜像切换，共用推理核心源码 &
两种系统分别启动，验证中断、线程切换、向量上下文和推理任务 \\
智能计算与视觉 &
统一算子描述、三后端执行和数值验证，首期使用受限小模型及测试图像输入 &
模型来源、预处理规则、逐层参考结果及板端结果对照 \\
感知与通信 &
预留张量输入、结果输出和通信适配接口，后续连接传感器、视觉设备与移动数据中心 &
采集设备及协议确认、驱动测试、真实输入与外部接口联调 \\
完整技术体系 &
衔接硬件参数、驱动、系统移植、模型转换与应用程序 &
版本记录、构建说明、接口约定和可重复的验证流程 \\
\bottomrule
\end{tabularx}
\end{table}

解题的关键在于保持各层约定一致。模型侧的形状、尺度与布局必须传递到算子调用，硬件侧的指令、地址图和加速器参数必须对应软件构建，操作系统侧则需要保证执行上下文、缓冲区和任务完成状态正确。三类后端采用同一数学定义，分别承担适合其结构的计算；后端选择不能改变模型的数值含义。\cite{technicalreport}

本阶段以现有试验箱和\TargetCost{}组织研究，不引入流片、自制产品板或采购平台。机器视觉和传感扩展保留明确的输入边界，“移动数据中心”的平台、协议和鉴权条件仍需与命题需求进一步对应。静态图像或串口测试数据用于验证推理链路，其证据范围与实时摄像头采集、传感设备接入及外部中心联调分别记录。

国产处理器要求是本方案尚待落实的一项命题条件。当前采用的海外开源处理器与加速器设计，以及 AMD/Xilinx FPGA 原型，不能证明已经使用国产 RISC-V 处理器芯片；开放指令架构、开源许可与芯片来源也不能相互替代。本阶段准确定位为基于开源 RISC-V 设计的 FPGA 原型及国产开源\RealTimeSystem{}适配研究，保留后续向国产处理器平台迁移的接口与研究空间。\cite{technicalreport}
